\documentclass[11pt,letterpaper]{article}

\usepackage[top=1in, bottom=1in, left=1in, right=1in, headheight=14pt]{geometry}
\usepackage{fontspec}
\setmainfont{DejaVu Serif}
\setsansfont{DejaVu Sans}
\setmonofont{DejaVu Sans Mono}

\usepackage{xcolor}
\usepackage{titlesec}
\usepackage{fancyhdr}
\usepackage{enumitem}
\usepackage{hyperref}
\usepackage{booktabs}
\usepackage{tabularx}
\usepackage{longtable}
\usepackage{colortbl}
\usepackage{multirow}
\usepackage{array}
\usepackage{parskip}
\usepackage{tcolorbox}
\usepackage{graphicx}
\usepackage{lastpage}

%% ─── COLOR SCHEME: Deep Space / OpenGL Viewport ───────────────────────────
\definecolor{primary}{HTML}{1A237E}       % Deep indigo — section headers
\definecolor{secondary}{HTML}{01579B}     % Electric blue — subsection headers
\definecolor{tertiary}{HTML}{4A148C}      % Shader purple — subsubsection headers
\definecolor{accent}{HTML}{0288D1}        % Bright blue — highlights, arrows
\definecolor{lightprimary}{HTML}{E8EAF6}  % Indigo tint — quote boxes
\definecolor{lightsecondary}{HTML}{E1F5FE}
\definecolor{lighttertiary}{HTML}{F3E5F5}
\definecolor{rowgray}{HTML}{F5F5F5}
\definecolor{bordergray}{HTML}{BDBDBD}
\definecolor{subtlegray}{HTML}{757575}
\definecolor{darktext}{HTML}{212121}
\definecolor{codegreen}{HTML}{1B5E20}
\definecolor{amber}{HTML}{E65100}

%% ─── SECTION FORMATTING ───────────────────────────────────────────────────
\titleformat{\section}
  {\Large\bfseries\sffamily\color{primary}}
  {\thesection}{1em}{}
  [\vspace{-0.5em}\textcolor{bordergray}{\rule{\textwidth}{0.4pt}}]

\titleformat{\subsection}
  {\large\bfseries\sffamily\color{secondary}}
  {\thesubsection}{1em}{}

\titleformat{\subsubsection}
  {\normalsize\bfseries\sffamily\color{tertiary}}
  {\thesubsubsection}{1em}{}

\titlespacing*{\section}{0pt}{1.5em}{0.8em}
\titlespacing*{\subsection}{0pt}{1.2em}{0.5em}
\titlespacing*{\subsubsection}{0pt}{0.8em}{0.3em}

%% ─── CUSTOM ENVIRONMENTS ──────────────────────────────────────────────────
\newcommand{\stageheader}[2]{%
  \noindent\colorbox{#2}{\makebox[\textwidth][l]{%
    \textcolor{white}{\bfseries\sffamily\hspace{6pt}#1\hspace{6pt}}}}%
  \vspace{0.5em}\par%
}

\newtcolorbox{asciibox}{
  colback=rowgray, colframe=bordergray,
  arc=1pt, boxrule=0.5pt,
  left=6pt, right=6pt, top=4pt, bottom=4pt,
  fontupper=\small\ttfamily,
  breakable
}

\newtcolorbox{quotebox}{
  colback=lightprimary, colframe=primary,
  arc=2pt, boxrule=0.5pt,
  left=12pt, right=12pt, top=8pt, bottom=8pt,
  breakable
}

\newtcolorbox{codebox}{
  colback=rowgray, colframe=secondary,
  arc=2pt, boxrule=0.5pt,
  left=8pt, right=8pt, top=6pt, bottom=6pt,
  fontupper=\footnotesize\ttfamily,
  breakable
}

\newcommand{\metafield}[2]{\textbf{\sffamily #1} #2\par}
\newcolumntype{L}[1]{>{\raggedright\arraybackslash}p{#1}}
\newcolumntype{C}[1]{>{\centering\arraybackslash}p{#1}}
\newcommand{\tableheader}[1]{\textbf{\sffamily\textcolor{white}{#1}}}

%% ─── HEADERS / FOOTERS ────────────────────────────────────────────────────
\pagestyle{fancy}
\fancyhf{}
\fancyhead[L]{\small\sffamily\textcolor{subtlegray}{NeHe OpenGL Intelligence Pack}}
\fancyhead[R]{\small\sffamily\textcolor{subtlegray}{GSD Mission Package}}
\fancyfoot[C]{\small\sffamily\textcolor{subtlegray}{Page \thepage\ of \pageref{LastPage}}}
\renewcommand{\headrulewidth}{0.4pt}
\renewcommand{\footrulewidth}{0pt}

%% ─── HYPERLINKS ───────────────────────────────────────────────────────────
\hypersetup{
  colorlinks=true,
  linkcolor=secondary,
  urlcolor=accent,
  pdfauthor={GSD Ecosystem / Tibsfox},
  pdftitle={NeHe OpenGL Intelligence Pack -- GSD Mission Package},
  pdfsubject={Vision to Mission Pipeline},
}

\begin{document}

%% ═══════════════════════════════════════════════════════════════════════════
%% TITLE PAGE
%% ═══════════════════════════════════════════════════════════════════════════
\thispagestyle{empty}
\begin{center}
\vspace*{2.5cm}

{\small\sffamily\textcolor{primary}{\textbf{GSD RESEARCH MISSION PACKAGE}}}

\vspace{0.3cm}
\textcolor{bordergray}{\rule{8cm}{0.4pt}}

\vspace{1.5cm}
{\fontsize{26}{32}\selectfont\bfseries\sffamily\textcolor{primary}{NeHe OpenGL\\[0.3em]Intelligence Pack}}

\vspace{0.8cm}
{\Large\sffamily\textcolor{secondary}{48 Lessons $\cdot$ All Examples $\cdot$ Upstream Integration}}

\vspace{0.5cm}
{\large\sffamily\itshape\textcolor{tertiary}{skill-creator Code, Documentation \& Working Examples}}

\vspace{2.5cm}
{\sffamily\textcolor{accent}{Vision $\rightarrow$ Research $\rightarrow$ Mission}}\\[0.3em]
{\small\sffamily\textcolor{accent}{Full Pipeline Execution}}

\vspace{2.5cm}
{\small\sffamily\textcolor{subtlegray}{Date: March 27, 2026 $\ |\ $ Status: Ready for GSD Orchestrator}}\\[0.3em]
{\small\sffamily\textcolor{subtlegray}{Activation Profile: Fleet (17 roles) $\ |\ $ Estimated: 8--12 context windows across 3 sessions}}\\[0.3em]
{\small\sffamily\textcolor{subtlegray}{Source: nehe.gamedev.net $\ |\ $ Archive: github.com/gamedev-net/nehe-opengl}}
\end{center}
\newpage

%% ─── TOC ──────────────────────────────────────────────────────────────────
\tableofcontents
\newpage

%% ═══════════════════════════════════════════════════════════════════════════
\stageheader{STAGE 1 --- VISION DOCUMENT}{primary}
%% ═══════════════════════════════════════════════════════════════════════════

\section{NeHe OpenGL Intelligence Pack --- Vision Guide}

\metafield{Date:}{March 27, 2026}
\metafield{Status:}{Cold-Start Mission / Research Compiled}
\metafield{Depends on:}{gsd-skill-creator (dev, v1.49+), gsd-upstream-intelligence-pack-v1\_43.md, gsd-silicon-layer-spec.md}
\metafield{Context:}{Adding the complete NeHe OpenGL tutorial corpus (48 lessons, 43 platform ports) to skill-creator as upstream intelligence, working code examples, and documentation --- the canonical 3D graphics pedagogy of the late 1990s to 2010s, preserved and modernized.}

\subsection{Vision}

There are corpora that defined entire generations of programmers. Jeff Molofee's NeHe Productions --- Neon Helium --- was that corpus for real-time 3D graphics. Between 1997 and 2012, the NeHe tutorials taught OpenGL to hundreds of thousands of developers: the ones who built the games, the simulations, the visualizers, the CAD tools. Every lesson was a handshake between a beginner and a coordinate system, between a curious programmer and a rendering pipeline.

The 48 lessons form a complete developmental arc: from the first blank window (Lesson 01) through texture mapping, lighting, blending, particle systems, shadow volumes, collision detection, and finally ArcBall rotation (Lesson 48). Each lesson builds on the last --- a curriculum so well-sequenced that it became the de facto OpenGL textbook for a decade. The GitHub archive preserves source code across 43 platform ports: C++, Python, Java, Delphi, Perl, Ruby, Scheme, Fortran, Assembly, and more.

The GSD ecosystem is about architectural intelligence: building systems that \textit{know} their domain deeply enough to adapt, explain, and generate within it. The NeHe corpus is exactly the kind of structured, layered knowledge that skill-creator should hold. Not as static reference, but as \textit{executable intelligence} --- a taxonomy of 3D concepts mapped to skills, with working modern implementations ready for agent deployment. When a future skill-creator session encounters a request involving vertex buffers or shadow stencils or particle systems, it should be able to reach into this corpus and surface the right lesson, the right pattern, the right code.

The modernization mandate is real: NeHe's original code runs on Windows with legacy Win32 API and deprecated glaux. The working examples for this mission target Python with ModernGL or PyOpenGL + GLFW --- cross-platform, pip-installable, runnable on any developer machine without a Visual C++ IDE from 2002.

\subsection{Problem Statement}

\begin{enumerate}
\item \textbf{No structured 3D graphics knowledge in skill-creator.} The current upstream intelligence architecture monitors Anthropic channels and skill specifications, but holds no domain knowledge about computer graphics, rendering pipelines, or mathematical foundations of 3D space. Requests involving OpenGL, shaders, or 3D math require the orchestrator to synthesize knowledge from scratch.

\item \textbf{Legacy code is non-runnable without archaeology.} The GitHub archive's C++ examples require Win32, glaux, and legacy build systems. Python ports exist but are fragmented across contributors with inconsistent quality. There is no single tested, modern, cross-platform implementation of all 48 lessons.

\item \textbf{No skill taxonomy maps graphics concepts to skill-creator primitives.} The NeHe lessons span categories (setup, primitives, transforms, textures, lighting, advanced effects, physics) that map naturally to skill clusters, but this mapping doesn't exist in machine-readable form. A skill that generates OpenGL code should know whether it's operating in the foundation tier or the expert tier.

\item \textbf{Upstream intelligence gap for graphics API evolution.} The transition from legacy OpenGL (fixed-function pipeline) to modern OpenGL (programmable shaders, VBOs, VAOs) is architecturally significant. Skill-creator has no model of this transition --- no awareness of deprecation boundaries, no ability to flag when generated code uses obsolete patterns.

\item \textbf{Missing DACP bundles for graphics skill handoffs.} Agent-to-agent handoffs in graphics tasks should pass structured three-part bundles (human intent + lesson reference + executable code), not ambiguous markdown. No DACP schema exists for OpenGL lesson delivery.
\end{enumerate}

\subsection{Core Concept}

\textbf{Model: Lesson as Skill Node, Corpus as Directed Graph}

Each NeHe lesson is a skill node with defined prerequisites, concepts introduced, and successors. The 48 lessons form a directed acyclic graph (DAG) --- the rendering knowledge graph --- that can be traversed by skill-creator to answer graphics questions, generate code at the right complexity level, or build multi-lesson learning sequences.

\subsubsection{Rendering Knowledge Graph --- Domain Map}

\begin{asciibox}
NEHE OPENGL CORPUS --- SKILL GRAPH OVERVIEW

TIER 0: FOUNDATION (L01-L05)
  L01 Window Setup ──► L02 Primitives ──► L03 Color ──► L04 Rotation ──► L05 3D Shapes

TIER 1: TEXTURE & LIGHT (L06-L12)
  L06 Texture Mapping ──► L07 Filters+Lighting+KB ──► L08 Blending
  L05 ──► L09 Moving Bitmaps ──► L10 3D World Loading
  L10 ──► L11 Flag/Wave Effect ──► L12 Display Lists

TIER 2: FONTS & EFFECTS (L13-L20)
  L12 ──► L13 Bitmap Fonts ──► L14 Outline Fonts ──► L15 Texture Mapped Fonts
  L08 ──► L16 Fog ──► L17 2D Texture Font ──► L18 Quadratics
  L08 ──► L19 Particle Engine ──► L20 Masking

TIER 3: ADVANCED RENDERING (L21-L28)
  L07 ──► L21 Lines+Antialiasing+Sound ──► L22 Bump+Multitex+Ext
  L06 ──► L23 Sphere Mapping ──► L24 Tokens+Scissor+TGA
  L25 Morphing+OBJ Load ──► L26 Stencil+Reflections ──► L27 Shadow Volumes
  L28 Glow+Motion Blur

TIER 4: SIMULATION & WORLDS (L29-L36)
  L29 Blitter+RAW+ASM ──► L30 BSP Collision ──► L31 Model Loading
  L32 Picking+Fogging+Multitex ──► L33 TGA Compression
  L34 Height Maps ──► L35 AVI Playback ──► L36 Radial Blur+RTT

TIER 5: EXPERT (L37-L48)
  L37 Cel-Shading ──► L38 Compressed Textures ──► L39 Physics Intro
  L39 ──► L40 Rope Physics ──► L41 ODE Virtual Worlds
  L42 Multiple Viewports ──► L43 FreeType Fonts
  L44 3D Lens Flare+Occlusion ──► L45 Vertex Buffer Objects (VBO)
  L46 Fullscreen FSAA ──► L47 CG Vertex Shader ──► L48 ArcBall Rotation

  CROSS-CUTTING:
  Math Foundation ──────────────────────────────────────► All tiers
  DACP Bundle Schema ───────────────────────────────────► All tiers
  Modern API Wrapper (PyOpenGL/ModernGL) ───────────────► All tiers
\end{asciibox}

\subsubsection{Module Architecture}

\begin{asciibox}
MODULE DECOMPOSITION (6 modules, Fleet activation)

MODULE 1: Taxonomy & Skill Mapping
  - Lesson-to-skill YAML taxonomy (48 entries)
  - Concept DAG in JSON-LD
  - DACP bundle schemas per lesson tier
  - skill-creator YAML for "opengl-lesson" skill type

MODULE 2: Foundation Track (L01-L12)
  - Modern Python implementations (PyOpenGL + GLFW)
  - Tested, cross-platform, pip-runnable
  - Inline documentation referencing NeHe concept

MODULE 3: Intermediate Track (L13-L24)
  - Font rendering via FreeType/Pygame
  - Effects: fog, particles, masking, bump mapping
  - Multi-texture and extension handling

MODULE 4: Advanced Track (L25-L36)
  - OBJ loading, stencil buffer, shadows
  - BSP collision, height maps, render-to-texture
  - TGA loading, AVI playback abstraction

MODULE 5: Expert Track (L37-L48)
  - Cel-shading, physics simulations (L39-L41)
  - VBOs, FSAA, CG shader translation
  - FreeType, ArcBall rotation math

MODULE 6: Upstream Intelligence Integration
  - opengl-intelligence.yaml for skill-creator
  - DACP bundle format for graphics handoffs
  - Deprecation map: legacy OpenGL → modern OpenGL
  - Documentation indexer for lesson lookup
\end{asciibox}

\subsection{Research Modules}

\subsubsection{Module 1 --- Taxonomy \& Skill Mapping}

Produces the machine-readable knowledge graph that skill-creator uses to reason about 3D graphics requests. Output is a YAML taxonomy file, a JSON-LD concept DAG, and a skill definition that registers the ``opengl-lesson'' skill type. This module is the intelligence backbone --- all other modules are data that flows through it.

\subsubsection{Module 2 --- Foundation Track (Lessons 01--12)}

Implements and tests the twelve foundation lessons in modern Python. These cover: OpenGL window setup, drawing triangles and quads, adding color, rotation, 3D shapes, texture mapping, lighting and keyboard input, blending, moving bitmaps in 3D, loading a 3D world, wave/flag effect, and display lists. Every implementation uses PyOpenGL + GLFW (cross-platform) or ModernGL where the lesson maps cleanly to a modern shader pattern. Tests verify visual output via framebuffer comparison.

\subsubsection{Module 3 --- Intermediate Track (Lessons 13--24)}

Implements lessons covering bitmap fonts, outline fonts, texture-mapped fonts, fog effects, 2D texture fonts, quadratics (spheres, cylinders, discs), particle engines, texture masking, lines and antialiasing with timing, bump mapping and multitexturing, sphere mapping, and token/extension detection with scissor testing. The font lessons require FreeType integration. Particle engine (L19) is particularly important as a skill pattern for procedural geometry generation.

\subsubsection{Module 4 --- Advanced Track (Lessons 25--36)}

Implements morphing between OBJ meshes, stencil buffer reflections and clipping, shadow volumes, glow and motion blur effects, BSP tree collision detection, model loading, picking with fog and multitexturing, TGA loading (compressed and uncompressed), height map landscape rendering, AVI playback abstraction (cross-platform video texture), and radial blur with render-to-texture. The BSP and shadow volume lessons are the most architecturally complex and require careful dependency management.

\subsubsection{Module 5 --- Expert Track (Lessons 37--48)}

Implements cel-shading (toon shader), compressed texture loading (DXT), introduction to physical simulations (spring-mass systems), rope physics (Verlet integration), virtual worlds with ODE physics, multiple viewport rendering, FreeType font rendering in OpenGL, 3D lens flare with occlusion testing, vertex buffer objects (the modern GPU data path), fullscreen antialiasing, CG vertex shader translation (to GLSL), and ArcBall rotation (the interaction pattern for 3D viewport navigation). This module requires the most adaptation from legacy to modern API.

\subsubsection{Module 6 --- Upstream Intelligence Integration}

Produces the artifact that makes all other modules useful inside skill-creator: the intelligence index that agents can query. Includes: the YAML taxonomy for skill registration, DACP three-part bundle schemas for graphics skill handoffs, a deprecation map that flags legacy OpenGL patterns and suggests modern equivalents, and a documentation index that agents can search by concept keyword, lesson number, or API function name.

\subsection{Chipset Configuration}

\begin{asciibox}
# nehe-intelligence-pack.chipset.yaml
# Fleet activation: 17 roles, 6 modules

agnus:
  role: Orchestration / DMA / Context Management
  model: claude-opus-4-6
  primary_tasks:
    - Taxonomy design (Module 1 schema architecture)
    - Cross-module DAG validation
    - DACP schema design
    - Legacy-to-modern API deprecation analysis
    - Cross-module pattern detection (ANALYST)
  token_budget: 45%

denise:
  role: Display / Output Rendering
  model: claude-sonnet-4-6
  primary_tasks:
    - Python code implementation (Modules 2-5)
    - Documentation generation
    - YAML taxonomy compilation
    - Test harness construction
    - Index page generation
  token_budget: 45%

paula:
  role: Audio / I/O / Anticipatory
  model: claude-haiku-4-5
  primary_tasks:
    - Shared schemas and templates
    - File scaffolding
    - Lesson metadata extraction
    - Test stub generation
  token_budget: 10%
\end{asciibox}

\subsection{Scope Boundaries}

\subsubsection{In Scope (v1.0)}

\begin{itemize}
\item All 48 NeHe lessons fully documented in skill-creator taxonomy
\item Working Python implementations for all 48 lessons (PyOpenGL + GLFW, cross-platform)
\item Automated test harness with framebuffer validation for lessons 01--20
\item Visual correctness spot-checks for lessons 21--48
\item DACP bundle schema for OpenGL skill handoffs
\item Deprecation map: legacy OpenGL (fixed-function) to modern OpenGL 3.3+ patterns
\item skill-creator YAML registration for ``opengl-lesson'' skill type
\item Upstream intelligence documentation index (searchable by concept, lesson, API)
\item README and integration guide for adding the pack to the dev branch
\end{itemize}

\subsubsection{Out of Scope}

\begin{itemize}
\item Non-Python platform ports (C++, Java, Delphi, etc.) --- defer to GitHub archive
\item Vulkan or Metal equivalents
\item WebGL/WebGPU translations
\item Real-time audio integration (L21 sound component --- stub only)
\item ODE physics library full integration (L41 --- interface documented, implementation deferred)
\item AVI video playback (L35 --- platform-specific; cross-platform stub provided)
\item Full GLSL shader replacements for all fixed-function pipeline patterns (deferred to v2.0)
\end{itemize}

\subsection{Success Criteria}

\begin{enumerate}
\item All 48 lessons are represented in the skill-creator taxonomy YAML with correct metadata (tier, prerequisites, concepts, API calls).
\item All 48 Python implementations run without error on Ubuntu 24 with PyOpenGL and GLFW installed via pip.
\item Foundation lessons (01--12) produce pixel-correct output matching reference screenshots (within 2\% RMSE on framebuffer comparison).
\item DACP bundle schemas validate against the three-part bundle specification (intent + data + code).
\item The deprecation map correctly flags at least 15 legacy OpenGL patterns with modern equivalents.
\item The documentation index returns the correct lesson for at least 30/30 spot-check concept queries.
\item The skill-creator YAML installs cleanly via the skill registration process on the dev branch.
\item All test harness tests pass in a fresh Ubuntu 24 environment (no undeclared dependencies).
\item Expert-tier lessons (37--48) include inline documentation explaining the mathematical concept, not just the code.
\item The ArcBall rotation implementation (L48) passes quaternion unit tests.
\item The VBO implementation (L45) uses only OpenGL 3.3 core profile features (no deprecated fixed-function calls).
\item The complete pack is deliverable as a single directory that can be dropped into the skill-creator dev branch.
\end{enumerate}

\subsection{Relationship to Other Documents}

\begin{center}
\rowcolors{2}{rowgray}{white}
\begin{tabularx}{\textwidth}{L{4.5cm}L{2.5cm}X}
\toprule
\rowcolor{primary}
\tableheader{Document} & \tableheader{Relationship} & \tableheader{Notes} \\
\midrule
gsd-upstream-intelligence-pack-v1\_43.md & Parent architecture & This pack extends the upstream intelligence layer into a new domain (3D graphics) using the Watch/Classify/Trace/Adapt/Brief cycle \\
gsd-skill-creator-analysis.md & Implementation host & The YAML taxonomy and skill definition integrate directly into skill-creator's six-step observe-detect-suggest loop \\
gsd-silicon-layer-spec.md & Execution model & Code examples follow the silicon layer's sandboxed execution patterns \\
gsd-instruction-set-architecture-vision.md & ISA mapping & OpenGL API calls can be modeled as a domain-specific instruction set within the GSD ISA framework \\
the-space-between.pdf & Mathematical spine & L-Systems, vector calculus, and complex plane concepts from the textbook underlie lessons 11, 25, 39--41 \\
\bottomrule
\end{tabularx}
\end{center}

\subsection{The Through-Line}

The NeHe tutorials are not just a programming curriculum. They are a record of how human beings learned to build worlds in silicon --- lesson by lesson, coordinate system by coordinate system, one glVertex call at a time. Jeff Molofee wrote them between 1997 and 2012, giving them away freely, and they became the substrate on which a generation of 3D programmers grew. The archive at \texttt{nehe.gamedev.net} is quiet now, but its lesson count --- 48 --- is almost exactly the number of NeHe's own life chapters worth remembering.

The GSD ecosystem is built on the Amiga Principle: remarkable outcomes through architectural elegance and specialized execution paths. The Amiga 500 --- 7 MHz, 512K RAM --- rendered animations that stunned engineers working with machines ten times more powerful, because its specialized chips (Agnus, Denise, Paula) each did one thing with surgical precision. The NeHe corpus teaches the same lesson in a different register: a single glTexCoord2f call, composed with a glVertex3f in a display list, inside a carefully managed state machine, produces effects that seem to exceed what the underlying math should allow.

Both are about the spaces between. The magic of a textured cube isn't in any individual vertex --- it's in the mapping between texture coordinate and geometry. The magic of a particle system isn't in any individual particle --- it's in the emergent field they form. This is the intelligence this pack adds to skill-creator: not a lookup table, but a \textit{compositional grammar} of 3D space.

\newpage

%% ═══════════════════════════════════════════════════════════════════════════
\stageheader{STAGE 2 --- RESEARCH REFERENCE}{secondary}
%% ═══════════════════════════════════════════════════════════════════════════

\section{NeHe OpenGL --- Research Reference}

\subsection{How to Use This Document}

Stage 2 provides the research findings that agents use to implement each module. Each subsection corresponds to a mission module. The lesson taxonomy (Section 2.2) is the primary data source for Module 1. Code patterns (Sections 2.3--2.7) inform the Python implementations. Section 2.8 covers the deprecation landscape. Section 2.9 is the source bibliography.

Key source repositories and references:
\begin{itemize}
\item \textbf{Primary archive:} \url{https://github.com/gamedev-net/nehe-opengl} (MIT license, 43 platform ports)
\item \textbf{Python port reference:} \texttt{/python/} directory in the archive (PyOpenGL + GLUT baseline)
\item \textbf{Linux/SDL port reference:} \texttt{/linuxsdl/} directory (clean C, cross-platform)
\item \textbf{Documentation PDF:} NeHe's OpenGL Tutorials (Andreas Lagotzki, 2004, 482 pages)
\item \textbf{Original site:} \url{http://nehe.gamedev.net} (1997--2012, Jeff Molofee ``NeHe'')
\end{itemize}

\subsection{Complete Lesson Taxonomy (All 48 Lessons)}

\begin{center}
\rowcolors{2}{rowgray}{white}
\begin{longtable}{C{1.2cm}L{4.8cm}L{2cm}L{1.4cm}X}
\toprule
\rowcolor{primary}
\tableheader{Lesson} & \tableheader{Title} & \tableheader{Tier} & \tableheader{Complexity} & \tableheader{Key Concepts} \\
\midrule
\endhead
01 & Setting Up An OpenGL Window & Foundation & Low & Context creation, rendering loop, GLFW \\
02 & Your First Polygon & Foundation & Low & glVertex, GL\_TRIANGLES, GL\_QUADS \\
03 & Adding Color & Foundation & Low & glColor3f, smooth shading, flat shading \\
04 & Rotation & Foundation & Low & glRotatef, animation loop, delta time \\
05 & 3D Shapes & Foundation & Low & glTranslatef, 3D primitives, depth buffer \\
06 & Texture Mapping & Foundation & Medium & glGenTextures, BMP loading, UV coordinates \\
07 & Filters, Lighting, Keyboard & Foundation & Medium & GL\_LIGHT0, glMaterialfv, key handling \\
08 & Blending & Foundation & Medium & glBlendFunc, GL\_SRC\_ALPHA, transparency \\
09 & Moving Bitmaps In 3D Space & Foundation & Medium & Billboard sprites, depth testing \\
10 & Loading A 3D World & Foundation & Medium & Text-based world format, sectors \\
11 & Flag Effect (Waving Texture) & Foundation & Medium & Sine-based vertex displacement, wave sim \\
12 & Display Lists & Foundation & Medium & glNewList, glCallList, GL\_COMPILE \\
13 & Bitmap Fonts & Intermediate & Medium & wglUseFontBitmaps, glBitmap, font metrics \\
14 & Outline Fonts & Intermediate & Medium & wglUseFontOutlines, 3D extruded glyphs \\
15 & Texture Mapped Outline Fonts & Intermediate & Medium & Font on rotating geometry \\
16 & Cool Looking Fog & Intermediate & Low & glFogf, GL\_EXP2, fog density \\
17 & 2D Texture Font & Intermediate & Medium & Texture atlas font, sprite sheets \\
18 & Quadratics & Intermediate & Medium & gluSphere, gluCylinder, gluDisk \\
19 & Particle Engine (Triangle Strips) & Intermediate & High & GL\_TRIANGLE\_STRIP, gravity, color fade \\
20 & Masking & Intermediate & Medium & Stencil-based masking, logical ops \\
21 & Lines, Antialiasing, Timing, Sound & Intermediate & High & GL\_LINE\_LOOP, GL\_SMOOTH, timer, audio \\
22 & Bump Mapping \& Multitexturing & Intermediate & High & GL\_TEXTURE1\_ARB, EMBM, extensions \\
23 & Sphere Mapping & Intermediate & Medium & GL\_SPHERE\_MAP, environment mapping \\
24 & Tokens, Scissor Testing, TGA & Intermediate & Medium & glScissor, TGA format, extension tokens \\
25 & Morphing \& OBJ Loading & Advanced & High & Vertex interpolation, Wavefront OBJ \\
26 & Stencil Buffer Reflections & Advanced & High & glStencilFunc, planar reflection \\
27 & Shadow Volumes & Advanced & High & Stencil shadow volumes, silhouette edges \\
28 & Glow \& Motion Blur & Advanced & High & Accumulation buffer, glow mask \\
29 & Blitter, RAW Textures, ASM & Advanced & High & Raw texture format, software blit \\
30 & Collision Detection (BSP Trees) & Advanced & Very High & BSP tree construction, AABB, ray cast \\
31 & Model Loading & Advanced & High & MD2 format, keyframe animation \\
32 & Picking, Fog, Multitexturing & Advanced & High & glSelectBuffer, render-for-pick \\
33 & Loading Compressed/Uncompressed TGA & Advanced & Medium & TGA RLE, BGRA conversion \\
34 & Height Maps & Advanced & High & Heightmap sampling, terrain mesh gen \\
35 & Playing AVI Files In OpenGL & Advanced & High & Video texture streaming, DirectShow \\
36 & Radial Blur, Render-to-Texture & Advanced & High & FBO, radial blur shader, RTT pattern \\
37 & Cel-Shading (Toon Shading) & Expert & High & 1D texture lookup, silhouette pass \\
38 & Loading Compressed Textures & Expert & Medium & DXT1/3/5, GL\_EXT\_texture\_compression \\
39 & Introduction To Physical Simulations & Expert & High & Spring-mass system, Euler integration \\
40 & Rope Physics & Expert & High & Verlet integration, constraint solving \\
41 & Virtual Worlds With ODE & Expert & Very High & Open Dynamics Engine, rigid body sim \\
42 & Multiple Viewports & Expert & Medium & glViewport, multiple cameras \\
43 & FreeType Fonts In OpenGL & Expert & High & FreeType2, glyph atlas, SDF readiness \\
44 & 3D Lens Flare With Occlusion & Expert & High & GL\_QUERY, occlusion test, flare fade \\
45 & Vertex Buffer Objects (VBO) & Expert & High & glGenBuffers, GL\_ARRAY\_BUFFER, VAO \\
46 & Fullscreen AntiAliasing (FSAA) & Expert & Medium & WGL\_SAMPLE\_BUFFERS, multisampling \\
47 & CG Vertex Shader & Expert & High & Cg toolkit, GLSL translation \\
48 & ArcBall Rotation & Expert & High & Quaternion math, trackball interaction \\
\bottomrule
\end{longtable}
\end{center}

\subsection{Module 2 Reference --- Foundation Track Implementation Patterns}

\subsubsection{Cross-Platform Window Setup (Lesson 01 --- Modern Python)}

The original NeHe Lesson 01 uses Win32 API (CreateWindowEx, WNDCLASS, wglCreateContext). The modern cross-platform equivalent uses GLFW via the \texttt{pyglfw} package. Key pattern:

\begin{codebox}
import glfw
from OpenGL.GL import *

def main():
    if not glfw.init():
        raise RuntimeError("GLFW init failed")
    glfw.window_hint(glfw.CONTEXT_VERSION_MAJOR, 3)
    glfw.window_hint(glfw.CONTEXT_VERSION_MINOR, 3)
    glfw.window_hint(glfw.OPENGL_PROFILE, glfw.OPENGL_CORE_PROFILE)
    window = glfw.create_window(800, 600, "NeHe L01 - OpenGL Window", None, None)
    glfw.make_context_current(window)
    glClearColor(0.0, 0.0, 0.0, 1.0)
    while not glfw.window_should_close(window):
        glClear(GL_COLOR_BUFFER_BIT | GL_DEPTH_BUFFER_BIT)
        glfw.swap_buffers(window)
        glfw.poll_events()
    glfw.terminate()
\end{codebox}

\textbf{Note on legacy vs. modern API:} Lessons 01--12 use the fixed-function pipeline (glBegin/glEnd, glRotatef, glTranslatef, glColor3f). The implementations preserve this for pedagogical fidelity --- the DACP skill metadata flags these as ``legacy-compatible'' and the deprecation map notes the modern VBO+shader equivalents.

\subsubsection{Texture Loading Pattern (Lesson 06)}

The BMP loading in the original uses Windows-specific bitmap functions. Modern replacement uses \texttt{Pillow} (PIL):

\begin{codebox}
from PIL import Image
import numpy as np

def load_texture(filename):
    img = Image.open(filename).convert("RGB")
    img_data = np.array(img, dtype=np.uint8)
    tex_id = glGenTextures(1)
    glBindTexture(GL_TEXTURE_2D, tex_id)
    glTexParameteri(GL_TEXTURE_2D, GL_TEXTURE_MIN_FILTER, GL_LINEAR)
    glTexParameteri(GL_TEXTURE_2D, GL_TEXTURE_MAG_FILTER, GL_LINEAR)
    glTexImage2D(GL_TEXTURE_2D, 0, GL_RGB, img.width, img.height,
                 0, GL_RGB, GL_UNSIGNED_BYTE, img_data)
    return tex_id
\end{codebox}

\subsection{Module 3 Reference --- Intermediate Track Key Patterns}

\subsubsection{Particle Engine (Lesson 19)}

The particle engine is one of the most important intermediate patterns --- it demonstrates the fundamental particle system architecture that appears in game engines, VFX tools, and GPU compute systems. Key data structure:

\begin{codebox}
import numpy as np
from dataclasses import dataclass

@dataclass
class Particle:
    active: bool = True
    life: float = 1.0
    fade: float = 0.025
    r: float = 1.0; g: float = 0.0; b: float = 0.0
    x: float = 0.0; y: float = 0.0; z: float = 0.0
    xi: float = 0.0; yi: float = 0.0; zi: float = 0.0
    xg: float = 0.0; yg: float = -0.8; zg: float = 0.0

MAX_PARTICLES = 1000
particles = [Particle(fade=np.random.uniform(0.003, 0.025))
             for _ in range(MAX_PARTICLES)]
\end{codebox}

\subsubsection{Bump Mapping / Multitexturing (Lesson 22)}

Requires OpenGL extensions detection. Modern pattern uses \texttt{OpenGL.extensions}:

\begin{codebox}
from OpenGL.extensions import hasGLExtension

if hasGLExtension('GL_ARB_multitexture'):
    from OpenGL.GL.ARB.multitexture import *
    glActiveTextureARB(GL_TEXTURE0_ARB)
    glBindTexture(GL_TEXTURE_2D, base_tex)
    glActiveTextureARB(GL_TEXTURE1_ARB)
    glBindTexture(GL_TEXTURE_2D, bump_tex)
\end{codebox}

\subsection{Module 4 Reference --- Advanced Track Key Patterns}

\subsubsection{Stencil Buffer Pattern (Lessons 26--27)}

The stencil buffer enables reflections and shadow volumes. Critical state management:

\begin{codebox}
# Reflection pattern (Lesson 26)
glEnable(GL_STENCIL_TEST)
glStencilFunc(GL_ALWAYS, 1, 0xFF)
glStencilOp(GL_KEEP, GL_KEEP, GL_REPLACE)
# Draw floor into stencil
glStencilFunc(GL_EQUAL, 1, 0xFF)
glStencilOp(GL_KEEP, GL_KEEP, GL_KEEP)
# Draw reflected scene (scaled -1 on Y)
glPushMatrix()
glScalef(1.0, -1.0, 1.0)
draw_scene()
glPopMatrix()
glDisable(GL_STENCIL_TEST)
\end{codebox}

\subsubsection{Height Map Terrain (Lesson 34)}

\begin{codebox}
def get_height(map_data, x, z, scale=10):
    # map_data: numpy array from grayscale image
    return float(map_data[z, x]) * scale / 255.0

def build_terrain(map_data, width, depth):
    vertices = []
    for z in range(depth - 1):
        for x in range(width):
            for dz in [0, 1]:
                h = get_height(map_data, x, z + dz)
                vertices.extend([x - width/2, h, z + dz - depth/2])
    return np.array(vertices, dtype=np.float32)
\end{codebox}

\subsection{Module 5 Reference --- Expert Track Key Patterns}

\subsubsection{Vertex Buffer Objects (Lesson 45)}

VBOs are the most architecturally significant modern OpenGL pattern in the NeHe series --- the transition from per-vertex CPU calls to GPU-resident data:

\begin{codebox}
# VBO + VAO pattern (Lesson 45 modernized)
import numpy as np
from OpenGL.GL import *

def create_vbo(vertices, indices):
    vao = glGenVertexArrays(1)
    vbo = glGenBuffers(1)
    ebo = glGenBuffers(1)
    glBindVertexArray(vao)
    glBindBuffer(GL_ARRAY_BUFFER, vbo)
    glBufferData(GL_ARRAY_BUFFER, vertices.nbytes, vertices, GL_STATIC_DRAW)
    glBindBuffer(GL_ELEMENT_ARRAY_BUFFER, ebo)
    glBufferData(GL_ELEMENT_ARRAY_BUFFER, indices.nbytes, indices, GL_STATIC_DRAW)
    glVertexAttribPointer(0, 3, GL_FLOAT, GL_FALSE, 12, ctypes.c_void_p(0))
    glEnableVertexAttribArray(0)
    glBindVertexArray(0)
    return vao, vbo, ebo
\end{codebox}

\subsubsection{ArcBall Rotation (Lesson 48)}

The ArcBall is the standard interaction model for 3D viewport navigation --- maps 2D mouse movement to 3D rotation via quaternions:

\begin{codebox}
import numpy as np
from numpy.linalg import norm

def map_to_sphere(point, width, height):
    """Map 2D screen point to unit sphere surface."""
    v = np.array([(2 * point[0] - width) / width,
                  (height - 2 * point[1]) / height, 0.0])
    d = norm(v[:2])
    v[2] = np.sqrt(max(1.0 - d * d, 0.0)) if d < 1.0 else 0.0
    return v / norm(v) if norm(v) > 0 else v

def compute_rotation(v1, v2):
    """Compute quaternion rotating v1 to v2 on unit sphere."""
    perp = np.cross(v1, v2)
    angle = np.dot(v1, v2)
    return np.array([perp[0], perp[1], perp[2], angle])  # xyzw quaternion
\end{codebox}

\subsubsection{FreeType Fonts (Lesson 43)}

\begin{codebox}
import freetype  # pip install freetype-py

def load_glyph_atlas(font_path, size=48):
    face = freetype.Face(font_path)
    face.set_pixel_sizes(0, size)
    characters = {}
    for c in range(128):
        face.load_char(chr(c), freetype.FT_LOAD_RENDER)
        g = face.glyph
        tex = glGenTextures(1)
        glBindTexture(GL_TEXTURE_2D, tex)
        glTexImage2D(GL_TEXTURE_2D, 0, GL_RED,
                     g.bitmap.width, g.bitmap.rows, 0,
                     GL_RED, GL_UNSIGNED_BYTE, g.bitmap.buffer)
        characters[chr(c)] = {
            'texture': tex, 'size': (g.bitmap.width, g.bitmap.rows),
            'bearing': (g.bitmap_left, g.bitmap_top),
            'advance': g.advance.x
        }
    return characters
\end{codebox}

\subsection{Module 6 Reference --- Upstream Intelligence Integration}

\subsubsection{DACP Bundle Schema for OpenGL Lessons}

Each OpenGL skill delivery follows the DACP three-part bundle structure:

\begin{codebox}
# DACP Bundle: OpenGL Lesson Handoff
# Part 1: Human Intent
intent:
  type: "opengl_lesson_request"
  lesson_id: 19
  tier: "intermediate"
  description: "Implement particle engine with triangle strips"
  target_language: "python"
  api_profile: "legacy_compatible"  # or "modern_core"

# Part 2: Structured Data
data:
  lesson_metadata:
    title: "Particle Engine Using Triangle Strips"
    prerequisites: [8, 9]  # lessons
    concepts: ["GL_TRIANGLE_STRIP", "particle_system", "gravity"]
    key_apis: ["glBegin", "GL_TRIANGLE_STRIP", "glColor4f"]
  deprecation_flags:
    - pattern: "glBegin/glEnd"
      severity: "info"
      modern_equivalent: "VBO + glDrawArrays"

# Part 3: Executable Code
code:
  entry_point: "lessons/lesson19/lesson19.py"
  dependencies: ["pyopengl", "pyglfw", "Pillow", "numpy"]
  test_command: "pytest tests/test_lesson19.py"
  run_command: "python lessons/lesson19/lesson19.py"
\end{codebox}

\subsubsection{Deprecation Map --- Legacy to Modern OpenGL}

\begin{center}
\rowcolors{2}{rowgray}{white}
\begin{tabularx}{\textwidth}{L{3.5cm}L{2.5cm}XL{1.8cm}}
\toprule
\rowcolor{primary}
\tableheader{Legacy Pattern} & \tableheader{OpenGL Version} & \tableheader{Modern Equivalent} & \tableheader{Severity} \\
\midrule
glBegin/glEnd & $\leq$2.1 & VBO + glDrawArrays / glDrawElements & Core \\
glRotatef/glTranslatef & $\leq$2.1 & GLM / numpy + uniform matrix & Core \\
gluPerspective & $\leq$2.1 & GLM perspective() or custom & Core \\
glVertex/glColor/glNormal & $\leq$2.1 & Vertex attributes in VBO & Core \\
GL\_LIGHTING state machine & $\leq$2.1 & Fragment shader & Core \\
Display lists (glNewList) & $\leq$2.1 & VAO + VBO + instancing & Core \\
gluSphere/gluCylinder & $\leq$2.1 & Custom geometry + shader & Info \\
Accumulation buffer & $\leq$2.1 & Multi-pass FBO + shader & Core \\
GLUT (glutMainLoop) & $\leq$2.1 & GLFW / SDL2 & Info \\
glaux library & $\leq$1.5 & Pillow / stb\_image & Core \\
glBitmap / wglUseFontBitmaps & Windows only & FreeType2 + texture atlas & Core \\
CG shader (cgc) & nVidia-only & GLSL 3.30+ & Core \\
GL\_QUADS primitive & $\leq$2.1 & Two GL\_TRIANGLES or GL\_TRIANGLE\_STRIP & Core \\
glTexEnvf (TexEnv combine) & $\leq$2.1 & Fragment shader texture sampling & Core \\
wgl* functions (context) & Windows only & GLFW / EGL / SDL2 context & Core \\
\bottomrule
\end{tabularx}
\end{center}

\subsection{Safety and Sensitivity Considerations}

\begin{center}
\rowcolors{2}{rowgray}{white}
\begin{tabularx}{\textwidth}{L{3.5cm}L{2cm}X}
\toprule
\rowcolor{primary}
\tableheader{Area} & \tableheader{Type} & \tableheader{Handling} \\
\midrule
Legacy code licensing & Copyright & All NeHe code is MIT licensed via the GitHub archive. All implementations must include the original copyright header and attribution to Jeff Molofee. \\
GPU driver variability & Technical & Tests must handle graceful fallback for missing extensions. FSAA (L46) and CG shaders (L47) require hardware-specific testing notes. \\
Physics simulation (L39-41) & Educational & Physics examples are demonstrations only. No real-world safety implications. ODE library dependency noted. \\
Platform-specific code (L35 AVI, L29 ASM) & Compatibility & Windows-specific components (DirectShow, x86 ASM) are replaced with cross-platform stubs. Stubs are clearly documented as non-functional placeholders pending platform-specific implementation. \\
\bottomrule
\end{tabularx}
\end{center}

\subsection{Source Bibliography}

\subsubsection{Primary Sources}

\begin{itemize}
\item Molofee, J. (``NeHe''), \textit{NeHe OpenGL Tutorials, Lessons 01--48}, Neon Helium Productions, 1997--2012. \url{http://nehe.gamedev.net}
\item gamedev-net (GitHub org), \textit{nehe-opengl: The complete archive of all NeHe OpenGL Lessons}, MIT License. \url{https://github.com/gamedev-net/nehe-opengl}
\item Lagotzki, A., \textit{NeHe's OpenGL Tutorials (PDF compilation, December 2004)}, 482 pages. Referenced as ``NeHe 2004 PDF.''
\end{itemize}

\subsubsection{Supporting Technical References}

\begin{itemize}
\item The Khronos Group, \textit{OpenGL 3.3 Core Profile Specification}. \url{https://registry.khronos.org/OpenGL/specs/gl/glspec33.core.pdf}
\item The Khronos Group, \textit{OpenGL Wiki --- Legacy OpenGL}. \url{https://www.khronos.org/opengl/wiki/Legacy_OpenGL}
\item PyOpenGL project, \textit{PyOpenGL 3.x Documentation}. \url{http://pyopengl.sourceforge.net/}
\item GLFW project, \textit{GLFW 3.x Documentation}. \url{https://www.glfw.org/docs/latest/}
\item FreeType project, \textit{FreeType 2 Documentation}. \url{https://freetype.org/freetype2/docs/}
\item ODE project, \textit{Open Dynamics Engine Documentation}. \url{https://www.ode.org/}
\item Shoemake, K. (1994), ``Arcball Rotation Control,'' \textit{Graphics Gems IV}, Academic Press. (ArcBall quaternion math reference for L48)
\end{itemize}

\newpage

%% ═══════════════════════════════════════════════════════════════════════════
\stageheader{STAGE 3 --- MISSION PACKAGE}{tertiary}
%% ═══════════════════════════════════════════════════════════════════════════

\section{Milestone Specification}

\subsection{Mission Objective}

Produce a fully integrated NeHe OpenGL Intelligence Pack for the gsd-skill-creator dev branch: a YAML taxonomy registering all 48 lessons as skill nodes, working Python implementations for all 48 lessons (PyOpenGL + GLFW, cross-platform, pip-installable), automated tests for lessons 01--20, DACP bundle schemas, a deprecation map from legacy to modern OpenGL, and a documentation index enabling agent-driven concept lookup --- delivered as a single droppable directory.

\subsection{Architecture Overview}

\begin{asciibox}
WAVE EXECUTION ARCHITECTURE

  ┌─────────────────────────────────────────────────────────┐
  │  WAVE 0: Foundation (Sequential, Paula/Haiku)            │
  │  Shared schemas · lesson metadata · dir scaffolding       │
  └──────────────────────────┬──────────────────────────────┘
                             │
         ┌───────────────────┼──────────────────────┐
         ▼                   ▼                      ▼
  ┌──────────────┐  ┌──────────────────┐  ┌─────────────────┐
  │ WAVE 1A      │  │  WAVE 1B         │  │  WAVE 1C        │
  │ Taxonomy &   │  │  Code: L01-L24   │  │  Code: L25-L48  │
  │ Intelligence │  │  Foundation +    │  │  Advanced +     │
  │ Integration  │  │  Intermediate    │  │  Expert         │
  │ (Agnus/Opus) │  │  (Denise/Sonnet) │  │  (Denise/Sonnet)│
  └──────┬───────┘  └────────┬─────────┘  └────────┬────────┘
         └───────────────────┼──────────────────────┘
                             │  CAPCOM GATE
                             ▼
  ┌─────────────────────────────────────────────────────────┐
  │  WAVE 2: Synthesis (Sequential, Agnus/Opus)              │
  │  Cross-module validation · DACP schemas · Deprecation map │
  │  Documentation index · Integration testing               │
  └──────────────────────────┬──────────────────────────────┘
                             │  CAPCOM GATE
                             ▼
  ┌─────────────────────────────────────────────────────────┐
  │  WAVE 3: Publication + Verification (Denise/Sonnet)      │
  │  README · skill-creator YAML · final test run · delivery  │
  └─────────────────────────────────────────────────────────┘
\end{asciibox}

\subsection{System Layers}

\begin{enumerate}
\item \textbf{Schema Layer} (Wave 0) --- Shared data structures: lesson metadata schema, DACP bundle schema, directory layout manifest
\item \textbf{Taxonomy Layer} (Wave 1A) --- YAML skill taxonomy, JSON-LD concept DAG, deprecation map, intelligence index
\item \textbf{Implementation Layer} (Waves 1B, 1C) --- 48 Python lesson files, organized by tier, with inline docs
\item \textbf{Test Layer} (Wave 1B/1C, finalized Wave 2) --- pytest suite for L01--L20, visual spot-check docs for L21--L48
\item \textbf{Integration Layer} (Wave 2) --- Cross-module linkage validation, DACP bundle generation, skill registration test
\item \textbf{Publication Layer} (Wave 3) --- README, CHANGELOG, dev branch integration guide, final delivery
\end{enumerate}

\subsection{Deliverables}

\begin{center}
\rowcolors{2}{rowgray}{white}
\begin{tabularx}{\textwidth}{C{0.5cm}L{4cm}XL{2cm}}
\toprule
\rowcolor{primary}
\tableheader{\#} & \tableheader{Deliverable} & \tableheader{Acceptance Criteria} & \tableheader{Spec} \\
\midrule
D1 & Lesson taxonomy YAML & 48 entries, valid schema, all fields populated & M1 \\
D2 & JSON-LD concept DAG & Valid JSON-LD, prerequisite edges correct & M1 \\
D3 & Deprecation map & 15+ patterns mapped, severity classified & M6 \\
D4 & DACP bundle schema & Validates against 3-part bundle spec & M6 \\
D5 & Documentation index & 30/30 concept queries return correct lesson & M6 \\
D6 & L01--L12 Python (×12) & Runs on Ubuntu 24, visual output correct & M2 \\
D7 & L13--L24 Python (×12) & Runs on Ubuntu 24, no dependency errors & M3 \\
D8 & L25--L36 Python (×12) & Runs on Ubuntu 24, stubs noted for L35 & M4 \\
D9 & L37--L48 Python (×12) & Runs on Ubuntu 24, L47 stub noted & M5 \\
D10 & Test suite (L01--L20) & All tests pass in fresh environment & M2/M3 \\
D11 & skill-creator YAML & Installs via skill registration process & M6 \\
D12 & README + integration guide & Single-command drop-in to dev branch & M6 \\
\bottomrule
\end{tabularx}
\end{center}

\subsection{Component Breakdown}

\begin{center}
\rowcolors{2}{rowgray}{white}
\begin{tabularx}{\textwidth}{L{3.5cm}L{3cm}L{1.8cm}L{2cm}X}
\toprule
\rowcolor{primary}
\tableheader{Component} & \tableheader{Dependencies} & \tableheader{Model} & \tableheader{Est. Tokens} & \tableheader{Notes} \\
\midrule
Lesson taxonomy YAML & lesson metadata schema & Haiku & 4k & Structured, templated \\
JSON-LD concept DAG & taxonomy YAML & Sonnet & 8k & Graph construction \\
Deprecation map & taxonomy, OpenGL spec & Opus & 12k & Requires API judgment \\
DACP bundle schema & 3-part bundle spec & Opus & 8k & Schema design \\
Documentation index & taxonomy + DAG & Sonnet & 6k & Index construction \\
L01--L05 Python & GLFW schema & Sonnet & 15k & Foundation, 5 files \\
L06--L12 Python & texture schema, light schema & Sonnet & 28k & Texture/light/blend \\
L13--L20 Python & font schema, particle schema & Sonnet & 32k & Complex patterns \\
L21--L28 Python & ext schema, stencil schema & Sonnet & 35k & Stencil/shadow heavy \\
L29--L36 Python & terrain schema, RTT schema & Sonnet & 38k & Simulation patterns \\
L37--L48 Python & physics schema, VBO schema & Opus+Sonnet & 55k & Expert tier, quaternions \\
Test suite L01-L20 & implementations, GL stub & Sonnet & 20k & pytest + framebuffer \\
skill-creator YAML & taxonomy, skill spec & Sonnet & 6k & Registration format \\
README + guide & all deliverables & Sonnet & 8k & Integration docs \\
\bottomrule
\end{tabularx}
\end{center}

\subsection{Model Assignment Rationale}

\begin{center}
\rowcolors{2}{rowgray}{white}
\begin{tabularx}{\textwidth}{L{2cm}C{1.8cm}XL{2.5cm}}
\toprule
\rowcolor{primary}
\tableheader{Model} & \tableheader{Budget \%} & \tableheader{Assigned Tasks} & \tableheader{Rationale} \\
\midrule
Opus & 30\% & Deprecation map design, DACP schema architecture, expert tier code (L37--L48), cross-module integration synthesis & Judgment-heavy: API evolution analysis, quaternion math, physics system design \\
Sonnet & 60\% & All foundation through advanced Python implementations (L01--L36), tests, taxonomy compilation, index construction, README & Structural implementation: high-volume, well-specified code generation \\
Haiku & 10\% & Lesson metadata schema, directory scaffolding, template stubs, test fixture generation & Fast scaffolding: templated, low-judgment work \\
\bottomrule
\end{tabularx}
\end{center}

\section{Wave Execution Plan}

\subsection{Wave Summary}

\begin{center}
\rowcolors{2}{rowgray}{white}
\begin{tabularx}{\textwidth}{C{1cm}L{3cm}L{1.8cm}L{2.5cm}X}
\toprule
\rowcolor{primary}
\tableheader{Wave} & \tableheader{Name} & \tableheader{Mode} & \tableheader{Model} & \tableheader{Output} \\
\midrule
0 & Foundation & Sequential & Haiku & Shared schemas, directory layout, lesson stubs \\
1A & Taxonomy Track & Parallel & Opus & YAML taxonomy, DAG, DACP schemas, deprecation map \\
1B & Foundation+Intermediate Code & Parallel & Sonnet & L01--L24 Python implementations + tests \\
1C & Advanced+Expert Code & Parallel & Sonnet+Opus & L25--L48 Python implementations + tests \\
GATE & CAPCOM Review & Sequential & Human & Scope confirm, visual checks, test results review \\
2 & Synthesis & Sequential & Opus & Cross-module validation, index, integration \\
GATE & CAPCOM Review & Sequential & Human & Final go/no-go before publication \\
3 & Publication & Sequential & Sonnet & README, skill YAML, final delivery \\
\bottomrule
\end{tabularx}
\end{center}

\subsection{Wave 0 --- Foundation}

\begin{center}
\rowcolors{2}{rowgray}{white}
\begin{tabularx}{\textwidth}{L{4cm}L{3cm}X}
\toprule
\rowcolor{primary}
\tableheader{Task} & \tableheader{Agent} & \tableheader{Output} \\
\midrule
Create directory structure & Paula (Haiku) & \texttt{nehe-intelligence-pack/\{lessons,taxonomy,tests,docs\}/} \\
Define lesson metadata schema & Paula (Haiku) & \texttt{schemas/lesson\_metadata.yaml} \\
Define DACP bundle template & Paula (Haiku) & \texttt{schemas/dacp\_bundle.yaml} \\
Generate lesson stub files & Paula (Haiku) & 48 × \texttt{lessons/lessonNN/lessonNN.py} (stubs) \\
Generate test fixture templates & Paula (Haiku) & \texttt{tests/conftest.py}, test stubs \\
\bottomrule
\end{tabularx}
\end{center}

\subsection{Wave 1 --- Parallel Tracks (A, B, C)}

\subsubsection{Track A --- Taxonomy and Intelligence}

\begin{center}
\rowcolors{2}{rowgray}{white}
\begin{tabularx}{\textwidth}{L{4cm}L{3cm}X}
\toprule
\rowcolor{primary}
\tableheader{Task} & \tableheader{Agent} & \tableheader{Output} \\
\midrule
Build YAML taxonomy (48 entries) & CRAFT-ARCH (Opus) & \texttt{taxonomy/nehe\_lessons.yaml} \\
Build JSON-LD concept DAG & CRAFT-ARCH (Opus) & \texttt{taxonomy/concept\_dag.json} \\
Design DACP bundle schemas & ARCH (Opus) & \texttt{schemas/dacp\_opengl.yaml} \\
Build deprecation map & ANALYST (Opus) & \texttt{taxonomy/deprecation\_map.yaml} \\
Design documentation index & INTEG (Sonnet) & \texttt{docs/index.yaml} \\
\bottomrule
\end{tabularx}
\end{center}

\subsubsection{Track B --- Foundation \& Intermediate Code (L01--L24)}

\begin{center}
\rowcolors{2}{rowgray}{white}
\begin{tabularx}{\textwidth}{L{4cm}L{3cm}X}
\toprule
\rowcolor{primary}
\tableheader{Task} & \tableheader{Agent} & \tableheader{Output} \\
\midrule
Implement L01--L05 (Foundation I) & EXEC-1 (Sonnet) & 5 × lesson py + tests \\
Implement L06--L12 (Foundation II) & EXEC-1 (Sonnet) & 7 × lesson py + tests \\
Implement L13--L18 (Intermediate I) & EXEC-2 (Sonnet) & 6 × lesson py + tests \\
Implement L19--L24 (Intermediate II) & EXEC-2 (Sonnet) & 6 × lesson py + tests \\
Run test suite L01--L20 & VERIFY (Sonnet) & Test results, defect log \\
Fix defects L01--L20 & EXEC-1/2 (Sonnet) & Patched implementations \\
\bottomrule
\end{tabularx}
\end{center}

\subsubsection{Track C --- Advanced \& Expert Code (L25--L48)}

\begin{center}
\rowcolors{2}{rowgray}{white}
\begin{tabularx}{\textwidth}{L{4cm}L{3cm}X}
\toprule
\rowcolor{primary}
\tableheader{Task} & \tableheader{Agent} & \tableheader{Output} \\
\midrule
Implement L25--L30 (Advanced I) & EXEC-3 (Sonnet) & 6 × lesson py \\
Implement L31--L36 (Advanced II) & EXEC-3 (Sonnet) & 6 × lesson py \\
Implement L37--L42 (Expert I) & CRAFT-IMPL (Opus+Sonnet) & 6 × lesson py + math docs \\
Implement L43--L48 (Expert II) & CRAFT-IMPL (Opus+Sonnet) & 6 × lesson py + math docs \\
Visual spot-check L21--L48 & VERIFY (Sonnet) & Screenshot docs + notes \\
\bottomrule
\end{tabularx}
\end{center}

\subsection{Cache Optimization Strategy}

\begin{itemize}
\item \textbf{Shared lesson schema:} Computed once in Wave 0, cached for all 48 implementation agents. Each agent loads the schema at start, reducing redundant schema reasoning.
\item \textbf{GLFW boilerplate:} L01 establishes the window + context pattern. All subsequent lessons inherit this via a shared \texttt{lessons/common/glfw\_window.py} module --- agents reference it, not re-implement it.
\item \textbf{Texture loading utility:} L06 establishes the Pillow-based texture loader. All texture-using lessons (L07--L48) import from \texttt{lessons/common/texture.py}.
\item \textbf{Taxonomy preload:} Track A output is cached at 5-minute TTL and preloaded into Track B/C context to enable DACP bundle population without re-computing metadata.
\item \textbf{Test fixtures:} Shared framebuffer comparison utility in \texttt{tests/conftest.py} loaded once per test session.
\end{itemize}

\subsection{Token Budget}

\begin{center}
\rowcolors{2}{rowgray}{white}
\begin{tabularx}{\textwidth}{L{3.5cm}C{2cm}C{2cm}C{2cm}X}
\toprule
\rowcolor{primary}
\tableheader{Wave/Track} & \tableheader{Opus (k)} & \tableheader{Sonnet (k)} & \tableheader{Haiku (k)} & \tableheader{Notes} \\
\midrule
Wave 0 & 0 & 0 & 20 & Scaffolding only \\
Wave 1A & 38 & 6 & 0 & Taxonomy + schemas \\
Wave 1B & 0 & 110 & 0 & L01--L24 code + tests \\
Wave 1C & 20 & 93 & 0 & L25--L48 code \\
Wave 2 & 25 & 20 & 0 & Synthesis + integration \\
Wave 3 & 0 & 22 & 5 & Publication \\
\midrule
\textbf{Total} & \textbf{83k} & \textbf{251k} & \textbf{25k} & $\approx$359k total \\
\bottomrule
\end{tabularx}
\end{center}

\subsection{Dependency Graph}

\begin{asciibox}
DEPENDENCY GRAPH

  [W0: Schemas]
       |
  ┌────┴──────────────────────────────────────────┐
  |                    |                          |
[W1A: Taxonomy]  [W1B: Code L01-24]  [W1C: Code L25-48]
  |                    |                          |
  |         [CAPCOM GATE: Visual check + test pass]
  |                    |                          |
  └────────────────────┴──────────────────────────┘
                       |
              [W2: Synthesis + Integration]
                  DACP bundles validated
                  Taxonomy cross-linked
                  Documentation index built
                       |
              [CAPCOM GATE: Final review]
                       |
              [W3: Publication]
                  skill-creator YAML
                  README + guide
                  Final delivery
\end{asciibox}

\section{Test Plan}

\subsection{Test Categories}

\begin{center}
\rowcolors{2}{rowgray}{white}
\begin{tabularx}{\textwidth}{L{3.5cm}C{1.5cm}L{2cm}X}
\toprule
\rowcolor{primary}
\tableheader{Category} & \tableheader{Count} & \tableheader{Priority} & \tableheader{Action on Fail} \\
\midrule
Safety-critical & 4 & P0 & BLOCK delivery \\
Core functionality & 48 & P1 & Fix before Wave 3 \\
Integration & 12 & P1 & Fix before Wave 3 \\
Regression & 20 & P2 & Document \& defer \\
Edge cases & 15 & P3 & Document \\
Visual spot-checks & 28 & P2 & Screenshot doc \\
\bottomrule
\end{tabularx}
\end{center}

\subsection{Safety-Critical Tests (BLOCK delivery if fail)}

\begin{center}
\rowcolors{2}{rowgray}{white}
\begin{tabularx}{\textwidth}{L{1.2cm}L{4cm}X}
\toprule
\rowcolor{primary}
\tableheader{ID} & \tableheader{Verifies} & \tableheader{Expected Result} \\
\midrule
SC-01 & No undeclared dependencies & \texttt{pip install -r requirements.txt} in fresh venv succeeds; all imports resolve \\
SC-02 & No Windows-only code in lessons 01--34 & \texttt{grep -r "win32\textbar{}glaux\textbar{}wgl"} returns zero matches outside explicitly documented stubs \\
SC-03 & L45 uses only OpenGL 3.3 core profile & No glBegin/glEnd, no fixed-function calls in lesson45.py; validated by code parser \\
SC-04 & All lesson files include MIT attribution & \texttt{grep -l "Jeff Molofee"} returns 48 files; validates copyright compliance \\
\bottomrule
\end{tabularx}
\end{center}

\subsection{Core Functionality Tests}

\begin{center}
\rowcolors{2}{rowgray}{white}
\begin{tabularx}{\textwidth}{L{2cm}L{3.5cm}X}
\toprule
\rowcolor{primary}
\tableheader{Test ID} & \tableheader{Lesson} & \tableheader{Test Method} \\
\midrule
CF-01 to CF-12 & L01--L12 & Headless framebuffer render (pytest-gl); compare against reference PNG within 2\% RMSE \\
CF-13 to CF-20 & L13--L20 & Import succeeds; init sequence completes; first frame renders without GL error \\
CF-21 & L45 VBO & Buffer binding, data upload, draw call all succeed; GL\_NO\_ERROR confirmed \\
CF-22 & L48 ArcBall & Unit tests on map\_to\_sphere() and compute\_rotation(); quaternion norm = 1.0 within 1e-6 \\
CF-23 & L43 FreeType & Glyph atlas generates 128 characters; no GL errors on atlas upload \\
CF-24 & L39--L40 Physics & Spring-mass and Verlet systems advance 1000 steps without NaN/Inf \\
\bottomrule
\end{tabularx}
\end{center}

\subsection{Verification Matrix}

\begin{center}
\rowcolors{2}{rowgray}{white}
\begin{tabularx}{\textwidth}{L{1cm}L{5cm}L{3.5cm}L{1.5cm}}
\toprule
\rowcolor{primary}
\tableheader{SC\#} & \tableheader{Success Criterion} & \tableheader{Test IDs} & \tableheader{Status} \\
\midrule
1 & All 48 lessons in taxonomy YAML & CF-T1 (schema validator) & Pre-Wave 1A \\
2 & All 48 Python implementations run & CF-01 to CF-20 + SC-01 & Wave 1B/1C \\
3 & Foundation visual correctness & CF-01 to CF-12 & Wave 1B \\
4 & DACP bundles validate & CF-T2 (DACP validator) & Wave 2 \\
5 & 15+ legacy patterns in deprecation map & SC-02 + manual count & Wave 1A \\
6 & Documentation index accuracy & CF-T3 (30 queries) & Wave 2 \\
7 & skill-creator YAML installs & CF-T4 (install test) & Wave 3 \\
8 & Tests pass in fresh environment & SC-01 & Wave 3 \\
9 & Expert lessons include math docs & Manual review (CF-T5) & Wave 1C \\
10 & ArcBall quaternion tests pass & CF-22 & Wave 1C \\
11 & L45 uses only 3.3 core profile & SC-03 & Wave 1C \\
12 & Single droppable directory & CF-T6 (structure check) & Wave 3 \\
\bottomrule
\end{tabularx}
\end{center}

\section{Mission Crew Manifest}

\begin{center}
\rowcolors{2}{rowgray}{white}
\begin{tabularx}{\textwidth}{L{2cm}L{3.5cm}L{2cm}X}
\toprule
\rowcolor{primary}
\tableheader{Role} & \tableheader{Agent} & \tableheader{Model} & \tableheader{Responsibility} \\
\midrule
FLIGHT & Orchestrator & Opus & Mission direction; wave go/no-go; scope arbitration \\
PLAN & Planner & Opus & Wave decomposition; dependency ordering; cache strategy \\
TOPO & Topology Manager & Sonnet & Agent team structure; module reassignment if needed \\
ARCH & Architecture & Opus & DACP schema design; deprecation map; taxonomy structure \\
SCOUT & Research Agent & Sonnet & Upstream source validation; PyOpenGL API research \\
EXEC-1 & Foundation Executor & Sonnet & L01--L12 implementation \\
EXEC-2 & Intermediate Executor & Sonnet & L13--L24 implementation \\
EXEC-3 & Advanced Executor & Sonnet & L25--L36 implementation \\
CRAFT-IMPL & Expert Implementer & Opus+Sonnet & L37--L48 implementation with math docs \\
CRAFT-ARCH & Taxonomy Architect & Opus & YAML taxonomy, JSON-LD DAG \\
ANALYST & Pattern Analyst & Opus & Cross-lesson pattern detection; deprecation analysis \\
VERIFY & Verification & Sonnet & Test execution; source validation; GL error checking \\
INTEG & Integration Agent & Sonnet & Cross-module linkage; index construction \\
SURGEON & Health Monitor & Sonnet & Code quality degradation detection \\
SECURE & Security & Sonnet & License compliance; copyright header validation \\
CAPCOM & HITL Interface & Human & Wave boundary approvals; visual checks \\
RETRO & Retrospective & Sonnet & Post-wave summaries; defect pattern analysis \\
LOG & Chronicle & Haiku & Commit messages; milestone summaries \\
\bottomrule
\end{tabularx}
\end{center}

\section{Execution Summary}

\begin{center}
\rowcolors{2}{rowgray}{white}
\begin{tabularx}{\textwidth}{L{5cm}X}
\toprule
\rowcolor{primary}
\tableheader{Metric} & \tableheader{Value} \\
\midrule
Total lessons covered & 48 \\
Platform ports in archive & 43 \\
Target implementation language & Python (PyOpenGL + GLFW) \\
Activation profile & Fleet (17 roles) \\
Estimated context windows & 8--12 \\
Estimated sessions & 3 (Wave 0+1 / Wave 2 / Wave 3) \\
Total token estimate & $\approx$359k \\
Opus \% & $\approx$23\% (83k) \\
Sonnet \% & $\approx$70\% (251k) \\
Haiku \% & $\approx$7\% (25k) \\
CAPCOM gates & 2 (post Wave 1; post Wave 2) \\
Safety-critical tests (BLOCK) & 4 \\
Core functionality tests & 24 \\
Total deliverables & 12 \\
Estimated Python lines of code & 6,000--9,000 (48 lessons × 125--190 LoC avg) \\
Dev branch integration & Single directory drop-in \\
License & MIT (Jeff Molofee / NeHe archive) \\
\bottomrule
\end{tabularx}
\end{center}

\vspace{2em}

\begin{quotebox}
\textit{``The Amiga rendered what IBM machines could only approximate --- not because it had more power, but because it had the right architecture. Agnus knew DMA. Denise knew color. Paula knew rhythm. Three chips, each doing one thing perfectly, composing effects that seemed to exceed what the underlying hardware should allow.''}

\smallskip

The NeHe corpus is the OpenGL equivalent of that moment: 48 lessons, each doing one thing perfectly, composing a complete education in how to build worlds from vertices and pixels. This intelligence pack doesn't just archive that knowledge --- it wires it into skill-creator's nervous system, so the lessons can be recalled, composed, and executed at machine speed, exactly when they're needed, in exactly the form they're needed in.

\smallskip

The spaces between the lessons --- the prerequisite edges in the DAG, the shared patterns that emerge across tiers, the moment a particle system and a physics simulator start to look like the same pattern at different scales --- that's where the intelligence lives.
\end{quotebox}

\end{document}
